\documentclass{amsart}
%\usepackage[sfdefault,lf]{carlito}
\usepackage[T1]{fontenc}
\renewcommand*\oldstylenums[1]{\carlitoOsF #1}
\usepackage[spanish]{babel}
\usepackage{titlesec}
\usepackage{hyperref}
\usepackage[dvipsnames]{xcolor} %https://www.overleaf.com/learn/latex/Using_colours_in_LaTeX
\usepackage{fancyhdr}
\usepackage{fontawesome5}%%http://mirrors.ibiblio.org/CTAN/fonts/fontawesome5/doc/fontawesome5.pdf
\usepackage{tikz}
\usepackage{geometry}
\usetikzlibrary{shapes, arrows.meta, positioning}
\pagestyle{fancy}
\definecolor{colorunison}{rgb}{0.69, 0.44, 0.09} %%% DEFINIMOS EL COLOR UNISON
\geometry{left=3cm, right=3cm,top=2cm}
\lhead{\textcolor{colorunison}{Título provisional del trabajo terminal }}
\rhead{\textcolor{colorunison}{\thepage}}
\cfoot{\textcolor{colorunison}{MCD UNISON}}

\titleformat{\section}
{\color{colorunison}\normalfont\Large\bfseries\sffamily}
{\color{colorunison}\thesection}{1em}{}

\titleformat{\subsection}
{\color{black!60}\normalfont\large\bfseries\sffamily}
{\color{black!60}\thesubsection}{1em}{}

\title{\sffamily Título provisional del trabajo terminal }
\author{Proyecto de trabajo terminal \\ Maestría en Ciencia de Datos}

\begin{document}



\begin{center}
\includegraphics[scale=0.3]{escudouson.png}
\end{center}

\maketitle

\thispagestyle{empty}

\section{Descripción General}

\noindent Queremos \textcolor{gray}{insertar el objetivo del proyecto} utilizando \textcolor{gray}{describir los datos principales a utilizar y su origen} para que \textcolor{gray}{describir en qué sentido va a cambiar el proyecto la forma de resolver un problema o cambios en la organización} y con esto \textcolor{gray}{describir el beneficio principal que se espera con el proyecto.}



\section{Información de contacto}

{\small 
\begin{enumerate}
\item[\faUserGraduate]  \textcolor{gray}{Nombre del estudiante}
\item[\faEnvelope] \textcolor{gray}{correo electrónico}
\end{enumerate}}
\subsubsection*{Director de trabajo terminal}
{\small
\begin{enumerate}
\item[\faUser]  \textcolor{gray}{Nombre del Director de trabajo terminal}
\item[\faEnvelope] \textcolor{gray}{correo electrónico}
\end{enumerate}}
\subsubsection*{Experto asesor \textcolor{gray}{(si aplica)}} 
{\small
\begin{enumerate}
\item[\faUser]  \textcolor{gray}{Nombre del experto asesor}
\item[\faBuilding ]  \textcolor{gray}{Nombre de la institución}
\item[\faEnvelope] \textcolor{gray}{correo electrónico}
\end{enumerate}}



\section{Propuesta} 
\noindent \textcolor{gray}{La propuesta debe responder las siguientes preguntas, algunas pueden omitirse si es que no se aplican al proyecto. Proponga subsecciones que respondan las siguientes preguntas, si  corresponde  (NO ESCRIBIR EN FORMATO PREGUNTA - RESPUESTA):  
\begin{itemize}
\item ¿Qué  problema  se  está  intentando  solucionar?
\item ¿Quiénes  se  verán  beneficiados  al  resolver  este problema?
\item¿Qué  se  está  haciendo  actualmente  para  solucionar  este  problema?
\item ¿Cómo  se  resolverá  el  problema  descrito  anteriormente?
\item ¿Por  qué  la  solución que  se propone  es  mejor?
\item ¿Cómo  se  define  el  éxito  de  esta  solución  y  cómo  se  evaluará  que  la  solución  efectivamente resuelve  la  problemática  planteada  anteriormente?
\item ¿Cómo  se  planea  garantizar  la  precisión  tecnológica  /  algorítmica  de  la  solución  propuesta  
y cómo  se  mantendrá  la  precisión  continua  de  esta  solución?
\item ¿Cómo  se  abordará  la  recopilación,  gestión  y  uso  ético  de  estos  datos,  incluida  la  
seguridad, privacidad,  integridad  y  transparencia  de  los  datos?
\item ¿Por  qué  tu  equipo  está  en  una  posición  única  para  resolver  este  problema  e  implementarla solución  propuesta?
\item ¿Además  de  los  miembros  listados  anteriormente,  quién  más  es  fundamental  para  el  éxito
de  este proyecto? ¿Cómo  se  van  a  construir  estas  estas  colaboraciones?
\end{itemize}}
\newpage


\thispagestyle{empty}

\noindent\textcolor{gray}{Ejemplo nombres de subsecciones:}

\



\subsection*{Planteamiento del problema}

\subsection*{Descripción de universo de  usuarios}

\subsection*{Soluciones existentes}  

\subsection*{Propuesta y justificación de solución}

\subsection*{Métricas de evaluación}

\subsection*{Posibilidades de escalamiento}

\subsection*{Descripción de conjunto de datos}

\subsection*{Propuesta de gestión de datos}

\subsection*{Equipo de trabajo}


\section{Cronograma de actividades}

\noindent \textcolor{gray}{Describir las actividades a realizar el tercer y el cuarto semestre.}

\section{Flujo de ciencia de datos (opcional)}

\noindent \textcolor{gray}{A  grandes  rasgos,  cuál  es  el  proceso  que  se  anticipa  seguir, esquema  para  diseñar,  implementar  y  evaluar  la  solución  propuesta.}
%\begin{center}
%    \begin{tikzpicture}[
%        node distance=4ex and 0em,
%        block/.style={rectangle, draw, fill=white!20, 
%    text width=5.5em, text centered, rounded corners, minimum height=4em},
%        line/.style={draw, -latex},
%        ]
%
%        \node [block] (1) {};
%        \node [block, below right= of 1] (2) {};
%        \node [block, below left= of 2] (3) {};
%        \node [block, above left= of 3] (4) {};
%
%        \path [line] (1.east) to[out=0, in=90] (2.60);
%        \path [line] (2.-60) to[out=-90, in=0] (3.east);
%        \path [line] (3.west) to[out=180, in=-90] (4.-120);
%        \path [line] (4.120) to[out=90, in=180] (1.west);
%    \end{tikzpicture}
% \end{center}











\end{document}